We report Monte Carlo experiments on rank-$r$ LoRA-like
ensembles validating the theoretical rates and constants. All
runs use the hard distribution $P^\star$ of \S\ref{sec:setup}
with task-dependent spectra $(D_t)_{t=1}^T$. Code, random seeds,
and raw outputs are released with the paper.

\subsection{Synthetic validation}\label{subsec:exp-synthetic}

\paragraph{Lower-bound slope.} We sample ensembles across
$T\in\{2,3,4,8\}$, $r\in\{4,8,16\}$, $d\in\{128,512,2048\}$,
with both isotropic spectra ($D_t = I_r$) and four
heavy-anisotropy families (geometric, linear,
two-scale, and mixed; spectra spread across
$[2^{-1}, 2^1]$). For each configuration we measure the excess
$\mathbb{E}_{P^\star}[\|\bar H^{1/2}w^\star - \eta\|^2]$ of the
Gaussian-QR + uniform scalar quantization algorithm of
\S\ref{sec:achv} as a function of $R = b\deff$ bits, and fit a
linear regression of $\log_2\mathbb{E}[\text{excess}]$ against
$b = R/\deff$. Across all tested configurations the empirical
slope is $-2.00 \pm 0.01$ (95\% bootstrap CI across 1000 trials
per configuration), matching the prediction $-2R/\deff$ of
Theorem~\ref{thm:general}. The floor term $B^2(1 - \E[\deff]/
(Tr))$ is reproduced to four to five digits of precision in
the common-$D_0$ and task-dependent-$D_t$ regimes.

\paragraph{Achievability constant.} We measure the ratio
$\mathbb{E}_{P^\star}[\max_t D_t(w^\star)] / (c_{\mathrm{TQ}}
B^2 2^{-2R/(Tr)})$ of empirical upper bound to the Stiefel
lower bound of Theorem~\ref{thm:general} at $c_{\mathrm{TQ}} = 1$.
The ratio is constant in $R$ — a necessary condition for the UB
to match the LB in order — and lies in $[11, 13]$ across
$T\in\{2,3\}$, $r = 4$, $d = 128$, for isotropic and
task-mismatched anisotropic spectra. This is consistent with
the analytical constant $C = Tc^2/3 \approx 17$ at $T=2$,
$c=5$, combined with $c_{\mathrm{TQ}} \approx 1$ from
TurboQuant~\cite{zandieh2025turboquant}.

\paragraph{Shared-$V$ regime.} For the floor-positive regime
we implement the null-space-aware bit allocation of
\S\ref{subsec:sharedv} with per-coord fractional bits and a
locked clip $c = 11.5\sigma_{\mathrm{pc}}$ (tuned once, applied
across all configurations). At $T = 2$, $r = 4$, across five
anisotropy mixes (iso+iso, iso+geom, geom+twoscale, lin+twoscale,
geom+iso), the empirical slope of excess over
$\mathrm{cheb}^2$ is $-1.60 \pm 0.10$ at $n_{\mathrm{trials}} =
1000$, matching the theoretical prediction
$-2r/(r + |A| - 1) = -8/5$ with bootstrap 95\% CI $\pm 0.010$
per cell. The $\pm 0.10$ spread reflects structural anisotropy:
iso+iso tracks closer to the scalar-quantization limit $-2$;
geom+twoscale undershoots due to $H_t$-metric distortion on the
perpendicular subspace. Extending to $T = 3, 4$ via a SOCP
Chebyshev solver with Gauss-Newton KKT refinement (residual
$\sim 10^{-15}$), the null-space split continues to track the
theoretical prediction across mixed anisotropies.

\paragraph{Random-codebook comparison.} As a sanity check on
the tightness of our linear construction in the shared-$V$
regime, we ran a valid rate-$R$ random-codebook encoder ($2^R$
iid Gaussian codewords, encoder picks the one minimizing
$\max_t D_t$, $R \leq 18$). At $T = 3$, $r = 3$, the
random-codebook slope is $-1.45$, strictly steeper than our
explicit linear construction's $-1.20$. This rules out
sharpening the lower bound to match the linear upper-bound
exponent; the true rate-distortion function in the
floor-positive regime lies strictly between $-2r/r = -2$ and
$-2r/(r + |A| - 1)$ and is left open
(Remark~\ref{rem:sharedv-gap}).

\subsection{Real-LLM merging}\label{subsec:exp-lora}

A calibration of Theorem~\ref{thm:general} against real LoRA
fine-tunes (measuring worst-task distortion on merged Qwen /
Gemma / Phi adapters for standard merging methods such as Task
Arithmetic, TIES, DARE, KnOTS, and TVQ) is deferred to a future
version of this paper. The aim of such experiments is to test
whether the bound predicts the \emph{ranking} of merging methods
across tasks of varying subspace overlap, not merely to
reproduce theoretical slopes on idealized data. We view this as
a substantive empirical contribution in its own right and
separate it from the present theoretical treatment.
